\section{Introduction}

Estimated time of arrival (ETA) is a central prediction task in transportation systems, informing navigation, dispatch, logistics, and traveler information services. In such settings, ETA is not merely a descriptive quantity: it shapes route choice, scheduling, service reliability, and user trust. For operational platforms, inaccurate or unstable ETA estimates propagate into missed coordination, inefficient routing, and degraded service quality~\cite{modica2024eta}.

The central difficulty addressed in this work is not the absence of traffic data in general, but the lack of open, reusable, route-aware data representations tailored to ETA-oriented tasks. Much of the traffic prediction literature relies on sensor-centric or aggregate spatio-temporal datasets, while trajectory-based data often captures movement without preserving the route-aware structure characteristic of navigation-style ETA settings~\cite{yu2017trafficpred}. This problem is compounded by the scarcity of openly reusable ETA datasets: many operationally valuable sources remain proprietary, making reproducible experimentation difficult. The gap is twofold: align road topology, traffic evolution, and route context in one coherent form, and support controlled downstream validation rather than stopping at offline export---ideally for both simulation-derived traffic and real trajectories in one framework.

This paper addresses that gap through a common ETA-oriented representation constructed from two fundamentally different data sources. The first path begins with controlled simulation, where the evolving traffic state and planned routes are directly observable. The second begins with real trajectories, where route inference and data repair are required before the data can be expressed in a coherent downstream form. Although these paths do not produce the same dataset, they produce distinct dataset instances expressed in the same graph-based representation. Roads are modeled as edges, while junctions and vehicles are modeled as nodes. The representation is explicitly route-aware: vehicles are tracked along known or inferred planned routes, making the data substantially closer to navigation-system usage than to generic traffic datasets. It is also dynamic in two senses: traffic evolution appears in time-varying node and edge features, and the graph structure itself changes as vehicles move and dynamic traffic relations are created and removed. While designed primarily for vehicle-level ETA-oriented tasks, the same representation can also support other downstream objectives, such as road-level load estimation or area-level traffic analysis.

The concrete contribution is an integrated artifact ecosystem composed of a desktop dataset-construction tool and a web-based testing environment.
The trajectory path is illustrated on the Porto taxi dataset (ECML/PKDD 2015 challenge; July 2013--June 2014)~\cite{porto2015pkdd} as a bundled conversion \emph{demonstrator}; other trajectories and geographies are supported. The desktop tool covers both paths: simulation scenarios with SUMO control, zone and fleet validation, and export of route-aware snapshots at chosen $\Delta t$; trajectories with network preparation, repair, route alignment, and export into the same representation. The web client attaches a compatible ETA predictor for interactive trips, logs outcomes to a database, and supports analysis, plots, and PDF reports; the open-source bundle allows substituting predictors and configurations. Artifact URLs are omitted for double-blind review and can be restored in the camera-ready version.

A bundled route-aware predictor exercises the toolchain end-to-end; Section~\ref{sec:eval} uses it only to demonstrate that the representation is learnable and usable in the web loop, not to advance predictive-model novelty.

The remainder of the paper positions the work (Section~\ref{sec:related}), presents the system overview (Section~\ref{sec:system}), the dataset-construction paths and shared representation (Section~\ref{sec:dataset}), the web-based evaluation loop (Section~\ref{sec:web}), and an evaluation snapshot with dataset statistics, metrics, and figures (Section~\ref{sec:eval}), followed by discussion and conclusion (Section~\ref{sec:discussion}).
